% Layout
% -------------------------------------------------------------------------
\usepackage[
    top=2cm,           % Увеличен отступ сверху, чтобы текст не налезал на колонтитул 
    bottom=1.5cm,      % Увеличен отступ снизу для корректного отображения и печати 
    left=1.5cm, 
    right=1.5cm,
    headheight=25pt,   % Фиксированная высота для устранения ворнингов fancyhdr 
    includehead        % Включает колонтитул в расчет области печати 
]{geometry}

\usepackage{fancyhdr}
\usepackage{multirow}
\usepackage{lastpage}
\usepackage{varwidth}

\pagestyle{fancy}
\fancyhf{} % Полная очистка текущих настроек 

% Настройки для одностороннего документа (как в вашем main.pdf) [cite: 2, 3]
\fancyhead[L]{\textit{\leftmark}} % Название главы слева 
\fancyhead[R]{\thepage/\pageref{LastPage}} % Номер страницы справа 

\fancyfoot[C]{} % Очистка центральной части футера 

% Стиль для страниц, где начинаются главы (обычно там пустые колонтитулы) 
\fancypagestyle{plain}{
    \fancyhf{}
    \renewcommand{\headrulewidth}{0pt}
}

\usepackage{titlesec}
\usepackage{ragged2e} % Пакет для правильного выравнивания без переносов

% Настройка формата главы
\titleformat{\chapter}[display]
  {\normalfont\huge\bfseries\RaggedRight} % Добавлено \RaggedRight для отмены переносов
  {\vspace{-2.5cm}\chaptertitlename\ \thechapter} 
  {10pt} 
  {\Huge}

% Настройка отступов главы
\titlespacing*{\chapter}{0pt}{0pt}{20pt}
% -------------------------------------------------------------------------